\appendix
\section{Supplementary Formal Material}
\label{sec:appendix}

This appendix collects formal content referenced from the main text but
not required to follow it: definitions that are not cited by number
elsewhere in the paper (Appendix~\ref{sec:appendix:definitions}), and the
standard theorems that follow from the formal model of
Section~\ref{sec:mathematical_model} but that this paper does not build
on directly (Appendix~\ref{sec:appendix:theorems}). The one theorem the
paper's argument actually depends on, cross-engine equivalence
(Theorem~\ref{thm:equivalence}), is proved in place in
Section~\ref{sec:correctness} rather than here.

\subsection{Additional Definitions}
\label{sec:appendix:definitions}

\begin{definition}[Entity]
\label{def:entity}
An \emph{entity} is a distinguished identifier $\varepsilon \in \mathcal{E} \subseteq \mathcal{I}$
with:
\begin{itemize}
    \item A unique identifier $\operatorname{id}(\varepsilon) \in \mathcal{I}$.
    \item A semantic type $\operatorname{type}(\varepsilon) \in \mathbf{SemanticType}$.
    \item An optional name $\operatorname{name}(\varepsilon)$.
    \item Optional attributes $\operatorname{attrs}(\varepsilon) \subseteq \mathcal{V}^*$.
\end{itemize}
If $f = (i, s, r, \vec{o}, c, k, m)$ is a fact, then $s \in \mathcal{E}$.
\end{definition}

\begin{definition}[Relation]
\label{def:relation}
A \emph{relation} is a distinguished identifier $\rho \in \mathcal{R} \subseteq \mathcal{I}$
with:
\begin{itemize}
    \item A unique identifier $\operatorname{id}(\rho) \in \mathcal{I}$.
    \item An arity $\operatorname{arity}(\rho) \geq 0$.
    \item Slot types $\operatorname{slots}(\rho) \in \mathbf{SemanticType}^*$.
\end{itemize}
A fact $f = (i, s, r, \vec{o}, c, k, m)$ is well-typed iff
$\operatorname{arity}(r) = |\vec{o}|$ and each $\omega_j$ conforms to
$\operatorname{slots}(r)_j$.
\end{definition}

\begin{definition}[Open Set]
\label{def:openset}
An \emph{open set} $U \in \mathcal{T}$ is a subset of $X$ that is declared open.
In the Alexandrov topology induced by $(\mathcal{C}, \leq)$, every down-set
$\downarrow c = \{d \in \mathcal{C} \mid d \leq c\}$ is open.
To each context $c$ we associate $U_c = \{f \in \mathcal{I} \mid \operatorname{context}(f) \leq c\}$.
Properties: $U_{\top} = X$; the assignment is monotone
($c_1 \leq c_2 \implies U_{c_1} \subseteq U_{c_2}$); intersections follow
Theorem~\ref{thm:meet} below ($U_{c_1 \wedge c_2}$ when the meet exists,
$\emptyset$ for incomparable contexts).
\end{definition}

\begin{definition}[Knowledge State]
\label{def:knowledgestate}
A \emph{knowledge state} is $\Sigma = (X, F)$ where $X$ is a finite topological space
of fact identifiers and $F$ is a presheaf on $X$ assigning facts to open sets.
State transitions: insert $\Sigma \mapsto \Sigma \cup \{f\}$,
delete $\Sigma \mapsto \Sigma \setminus \{f\}$,
update $\Sigma \mapsto (\Sigma \setminus \{f\}) \cup \{f'\}$.
\end{definition}

\begin{definition}[Canonical Query]
\label{def:canonicalquery}
A \emph{canonical query} is $Q = (\tau, s, r, \vec{o}, c, k_{\min}, n_{\max})$ where:
$\tau \in \{\text{FACT}, \text{WALK}, \text{JOIN}, \text{GLUE}\}$ is the query mode,
$s \in \mathcal{E} \cup \{\ast\}$ is the subject (wildcard $\ast$ matches all),
$r \in \mathcal{R} \cup \{\ast\}$ is the relation,
$\vec{o} \in \mathcal{V}^* \cup \{\ast\}$ is the object pattern,
$c \in \mathcal{C} \cup \{\ast\}$ is the context,
$k_{\min} \in [0,1]$ is the minimum confidence,
$n_{\max} \in \mathbb{N}$ is the result limit.
The denotation $\llbracket Q \rrbracket_\Sigma$ is the set of facts matching $Q$.
\end{definition}

A \emph{canonical result} is $R = \llbracket Q \rrbracket_\Sigma$, a set of
canonical facts. Section~\ref{sec:correctness} defines \emph{semantic
equivalence} on canonical results (Definition~\ref{def:semanticequiv}) as
literal equality of order-independent, type-preserving canonical result
sets, which is what the benchmark harness checks and what
Theorem~\ref{thm:equivalence} proves; that is the only definition of
result equivalence used anywhere in this paper.

\subsection{Theorems and Proofs}
\label{sec:appendix:theorems}

The definitions above and in Section~\ref{sec:formal_definitions} satisfy
the standard consequences one would expect of a presheaf formalisation.
None of the theorems below is cited by number elsewhere in the paper; we
record them here as a completeness check on the formal model rather than
as load-bearing results, in contrast to Theorem~\ref{thm:equivalence}
(Section~\ref{sec:correctness}), which the paper's central claim actually
depends on.

\subsubsection{Semantic Preservation}
\label{sec:thm:preservation}

\begin{theorem}[Semantic Preservation / Round-trip Correctness]
\label{thm:preservation}
Let $T: \mathcal{F} \to \mathcal{T}(\mathcal{F})$ be the triple decomposition
function (reifying an $n$-ary fact into $2+n$ triples), and
$R: \mathcal{T}(\mathcal{F}) \to \mathcal{F}$ be the reconstruction function.
For any semantic fact $f \in \mathcal{F}$:
\[
R(T(f)) = f.
\]
\end{theorem}

\begin{proof}[Proof Sketch]
Decomposition $T$ creates triples $(s, \mathtt{rdf:type}, \mathtt{rdf:Statement})$,
$(s, \mathtt{rdf:subject}, f.\text{subject})$, $(s, \mathtt{rdf:predicate}, f.\text{relation})$,
and $(s, \mathtt{rdf:object}_j, f.\text{objects}_j)$ for each $j$.
Reconstruction $R$ groups triples by statement ID $s$ and inverts each step.
Since $T$ is injective and $R$ is its left inverse on $\operatorname{im}(T)$,
the round-trip is identity.
\end{proof}

\subsubsection{Functoriality of Restriction}

\begin{theorem}[Restriction Functoriality]
\label{thm:restriction}
The restriction maps $\rho_{d,c}: F(d) \to F(c)$ for $c \leq d$ satisfy:
\[
\rho_{c,c} = \operatorname{id}_{F(c)}, \qquad
\rho_{e,c} = \rho_{d,c} \circ \rho_{e,d} \quad (c \leq d \leq e).
\]
Hence $F: \mathcal{C}^{\mathrm{op}} \to \mathbf{Set}$ is a functor (a presheaf).
\end{theorem}

\subsubsection{Sheaf Axioms}

Locality and gluing are \emph{axioms} that distinguish sheaves among
presheaves; for an arbitrary presheaf on an arbitrary space, both fail
in general. We record their statements here in the specific setting this
paper uses --- a presheaf induced by a functor on the finite context
poset, over the Alexandrov topology --- where, as discussed in
Section~\ref{sec:sheaf-remark}, every open set is a union of minimal
open neighbourhoods and both conditions are satisfied automatically.
They are stated for completeness, not as a technical contribution.

\begin{theorem}[Locality, Alexandrov setting]
\label{thm:locality}
Let $F$ be the sheaf on $(X, \mathcal{T}_A)$ induced by a functor on the
context poset, $U \in \mathcal{T}_A$, $\{U_i\}$ an open cover of $U$,
and $s, t \in F(U)$.
If $\rho_{U,U_i}(s) = \rho_{U,U_i}(t)$ for all $i$, then $s = t$.
\end{theorem}

\begin{proof}[Proof Sketch]
Every $x \in U$ lies in some $U_i$, and $\mathcal{B}(x) \subseteq U_i$,
so $s$ and $t$ agree on every minimal neighbourhood contained in $U$;
sections of the induced sheaf are determined by their values on minimal
neighbourhoods.
\end{proof}

\begin{theorem}[Gluing, Alexandrov setting]
\label{thm:gluing}
Let $F$ be as above, $U \in \mathcal{T}_A$,
$\{U_i\}$ an open cover of $U$, and $\{s_i \in F(U_i)\}$ a compatible family:
$\rho_{U_i, U_i \cap U_j}(s_i) = \rho_{U_j, U_i \cap U_j}(s_j)$ for all $i,j$.
Then there exists a unique $s \in F(U)$ with $\rho_{U,U_i}(s) = s_i$ for all $i$.
\end{theorem}

\begin{proof}[Proof Sketch]
A section of the induced sheaf over $U$ is an assignment of a value to
each minimal neighbourhood $\mathcal{B}(x) \subseteq U$, compatible
under restriction. For each $x \in U$, pick $i$ with $x \in U_i$ and
take the value $s_i$ assigns to $\mathcal{B}(x)$; compatibility ensures
this is well-defined on overlaps. Uniqueness follows from
Theorem~\ref{thm:locality}.
\end{proof}

\subsubsection{Canonical Equivalence}

\begin{theorem}[Invertibility of the Canonical Mapping]
\label{thm:canonical}
Let $\phi: \mathcal{F} \to \mathfrak{F}$ be the canonical mapping from facts to canonical facts,
and $\psi: \mathfrak{F} \to \mathcal{F}$ the inverse mapping. Then:
\[
\psi(\phi(f)) = f,\quad \phi(\psi(\mathfrak{f})) = \mathfrak{f}
\]
for all $f \in \mathcal{F}$, $\mathfrak{f} \in \mathfrak{F}$.
\end{theorem}

Cross-engine query equivalence between $\Sigma_{\mathrm{KG}}$ and
$\Sigma_{\mathrm{Sheaf}}$ is stated and proved as
Theorem~\ref{thm:equivalence} in Section~\ref{sec:correctness}, together
with the canonical-form machinery (Definitions~\ref{def:lossless}
and~\ref{def:semanticequiv}) it depends on; we do not restate it here to
avoid a second, weaker copy of the same claim.

\subsubsection{Structural Properties}

\begin{theorem}[Open-Set Intersections]
\label{thm:meet}
For any $c_1, c_2 \in \mathcal{C}$ in the Alexandrov topology $\mathcal{T}_A$:
\[
U_{c_1} \cap U_{c_2} =
\begin{cases}
U_{c_1 \wedge c_2} & \text{if } c_1, c_2 \text{ are comparable (the meet
is the more specific of the two)},\\
\emptyset & \text{otherwise.}
\end{cases}
\]
\end{theorem}

\begin{proof}
A fact identifier lies in $U_{c_1} \cap U_{c_2}$ iff its context $d$
satisfies $d \leq c_1$ and $d \leq c_2$. If $c_1 \leq c_2$, the two
conditions reduce to $d \leq c_1 = c_1 \wedge c_2$ (symmetrically for
$c_2 \leq c_1$). If $c_1$ and $c_2$ are incomparable, then $d \leq c_1$
and $d \leq c_2$ would require both $c_1$ and $c_2$ to be prefixes of
$d$, which is impossible: the two paths disagree at their first differing
segment. Hence no such $d$ exists and the intersection is empty.
\end{proof}

Note that the longest common prefix of two incomparable contexts is their
\emph{join} $c_1 \vee c_2$ (most specific common ancestor), not their
meet, and it bounds the union rather than the intersection:
$U_{c_1} \cup U_{c_2} \subseteq U_{c_1 \vee c_2}$, with equality failing
in general (the join's open set also contains facts from sibling
branches).

\begin{theorem}[Preservation of Referential Integrity]
\label{thm:integrity}
For any knowledge state $\Sigma$ and fact $f \in \Sigma$:
\[
\operatorname{subject}(f) \in \mathcal{E}_\Sigma,\quad
\operatorname{relation}(f) \in \mathcal{R}_\Sigma,\quad
\forall \omega_j \in \vec{o}: \omega_j \in \mathcal{E}_\Sigma \text{ if reference}.
\]
\end{theorem}

\begin{theorem}[Determinism of Query Execution]
\label{thm:determinism}
For a fixed query $Q$ and fixed knowledge state $\Sigma$,
the result $\llbracket Q \rrbracket_\Sigma$ is uniquely determined,
independent of execution order, caching state, parallelism, or query history.
\end{theorem}

\begin{theorem}[Consistency of Global Sections]
\label{thm:global}
A family of local sections $\{s_i \in F(U_i)\}$ can be glued to a global section
$s \in F(\mathcal{C})$ iff for every pair $(i,j)$:
\[
\rho_{U_i, U_i \cap U_j}(s_i) = \rho_{U_j, U_i \cap U_j}(s_j).
\]
\end{theorem}

\subsection{Further Detail}
\label{sec:appendix:further}

Fully worked proofs of every theorem above, pseudocode for the
algorithms described in Section~\ref{sec:system-architecture}, and the
implementation-to-definition validation table are kept in the repository
rather than reproduced here, since they change with the code and would
go stale in a typeset PDF:

\begin{description}
    \item[Proofs] \texttt{research/proofs/}
    \item[Algorithms] \texttt{research/mathematics/algorithms.md}
    \item[Validation table] \texttt{research/mathematics/validation.md}
\end{description}
